\clearpage
\SourceLeaf{607}
\NoSourceHead
\begingroup
\newcommand{\ContentsRow}[5]{\par\noindent\hangindent=#1\hspace*{#1}#2\hspace{.4em}\mbox{\footnotesize #3}\hfill\mbox{#4\quad #5}\par}
\begin{center}\Large\bfseries\textls[120]{INHALT.}\end{center}
\begin{center}\large\bfseries GAUSS WERKE BAND X{\normalsize 1}.\\NACHTRÄGE ZUR REINEN MATHEMATIK.\end{center}
\begin{center}\rule{.1\linewidth}{.4pt}\end{center}
\begin{center}\bfseries KLEINERE VERÖFFENTLICHUNGEN.\end{center}
\ContentsRow{0pt}{Zur Kreisteilung}{.\,.\,.\,.\,.\,.\,.\,.\,.\,.\,.\,.\,.\,.\,.\,.\,.\,.\,.\,.\,.\,.\,.\,.\,.\,.\,.\,.\,.}{Seite}{3}
\ContentsRow{1.4em}{Bemerkung}{.\,.\,.\,.\,.\,.\,.\,.\,.\,.\,.\,.\,.\,.\,.\,.\,.\,.\,.\,.\,.\,.\,.\,.\,.\,.\,.\,.\,.}{—}{4}
\ContentsRow{0pt}{Vorrede zum \emph{Lehrbuch der Astronomie} von \textsc{Joseph Piazzi}, übersetzt von J.\ H.\ \textsc{Westphal}}{.\,.}{—}{5}
\ContentsRow{0pt}{Eine in Deutschland erfundene Rechenmaschine}{.\,.\,.\,.\,.\,.\,.\,.\,.\,.\,.\,.\,.\,.\,.\,.\,.\,.}{—}{6}
\ContentsRow{1.4em}{Bemerkung}{.\,.\,.\,.\,.\,.\,.\,.\,.\,.\,.\,.\,.\,.\,.\,.\,.\,.\,.\,.\,.\,.\,.\,.\,.\,.\,.\,.\,.}{—}{6}
\ContentsRow{0pt}{Vorrede zu den \emph{Mathematischen Abhandlungen} von G.\ \textsc{Eisenstein}}{.\,.\,.\,.\,.\,.\,.\,.\,.\,.\,.}{—}{7}
\begin{center}\rule{.13\linewidth}{.4pt}\end{center}
\begin{center}\bfseries ARITHMETIK.\end{center}
\noindent\emph{Nachlass und Briefwechsel.}\par
\ContentsRow{.5em}{Einige asymptotische Gesetze der Zahlentheorie}{.\,.\,.\,.\,.\,.\,.\,.\,.\,.\,.\,.\,.\,.\,.\,.\,.}{—}{11}
\ContentsRow{1.9em}{Bemerkungen}{.\,.\,.\,.\,.\,.\,.\,.\,.\,.\,.\,.\,.\,.\,.\,.\,.\,.\,.\,.\,.\,.\,.\,.\,.\,.\,.}{—}{16}
\ContentsRow{.5em}{Zur Entstehungsgeschichte der Disquisitiones arithmeticae. \textsc{Gauss} an v.\ \textsc{Zimmermann} 1797 März 12}{.\,.\,.\,.\,.\,.\,.\,.\,.\,.\,.\,.\,.\,.\,.\,.\,.\,.\,.\,.\,.\,.\,.\,.\,.\,.\,.\,.\,.\,.}{—}{19}
\ContentsRow{1.9em}{Bemerkung}{.\,.\,.\,.\,.\,.\,.\,.\,.\,.\,.\,.\,.\,.\,.\,.\,.\,.\,.\,.\,.\,.\,.\,.\,.\,.\,.\,.}{—}{21}
\par\medskip\noindent\textls[120]{Zur Theorie der Potenzreste.}\par
\noindent\hspace*{.5em}\textls[120]{Quadratische Reste.}\par
\ContentsRow{1em}{I.\quad Über einige Summen}{.\,.\,.\,.\,.\,.\,.\,.\,.\,.\,.\,.\,.\,.\,.\,.\,.\,.\,.\,.\,.\,.\,.\,.\,.}{—}{22}
\ContentsRow{1.9em}{\textsc{Gauss} an \textsc{Olbers}, 1805 September 3}{.\,.\,.\,.\,.\,.\,.\,.\,.\,.\,.\,.\,.\,.\,.\,.\,.\,.}{—}{24}
\ContentsRow{3.3em}{Bemerkung}{.\,.\,.\,.\,.\,.\,.\,.\,.\,.\,.\,.\,.\,.\,.\,.\,.\,.\,.\,.\,.\,.\,.\,.\,.\,.}{—}{25}
\ContentsRow{1em}{II.\quad Bestimmung des Dezidenten in der Lehre von den quadratischen Resten}{.\,.\,.\,.\,.\,.}{—}{26}
\ContentsRow{3.3em}{Bemerkung}{.\,.\,.\,.\,.\,.\,.\,.\,.\,.\,.\,.\,.\,.\,.\,.\,.\,.\,.\,.\,.\,.\,.\,.\,.\,.}{—}{27}
\ContentsRow{1em}{II.\quad Fünfter Beweis des Fundamentalsatzes bei den quadratischen Resten}{.\,.\,.\,.\,.\,.\,.}{—}{28}
\ContentsRow{3.3em}{Bemerkung}{.\,.\,.\,.\,.\,.\,.\,.\,.\,.\,.\,.\,.\,.\,.\,.\,.\,.\,.\,.\,.\,.\,.\,.\,.\,.}{—}{32}
\endgroup
